% from ega4-2-fr.tex
\oldpage[IV-2]{104}

donc $\mathscr F$ est encore un $\mathcal O_X$-Module de Cohen-Macaulay en
ces points. Dire que $\mathscr F$ possède la propriété $(S_k)$ pour tout
$k$ signifie donc que $\mathscr F$ est un $\mathcal O_X$-Module de
Cohen-Macaulay. Il est clair que pour $k'\geq k$, la propriété $(S_{k'})$
implique $(S_k)$; pour $k\leq0$, tout $\mathcal O_X$-Module cohérent a la
propriété $(S_k)$.

\item[(ii)] Pour vérifier la condition (5.7.2.1), on peut se borner au cas
où $x\in\operatorname{Supp}(\mathscr F)$; dans le cas contraire on a en
effet $\dim(\mathscr F_x)=-\infty$ (0, 14.1.2).

\item[(iii)] Si $X=\operatorname{Spec}(A)$, où $A$ est un anneau
noethérien, et $\mathscr F=\widetilde M$, où $M$ est un $A$-module de type
fini, on dit que $M$ possède la propriété $(S_k)$ si $\mathscr F$ possède
cette propriété. Pour un préschéma localement noethérien quelconque $X$,
dire que $\mathscr F$ possède la propriété $(S_k)$ en un point $x\in X$
signifie donc que si on pose $Y=\operatorname{Spec}(\mathcal O_x)$, le
$\mathcal O_Y$-Module $\widetilde{\mathscr F_x}$ possède la propriété
$(S_k)$; on notera que la condition (5.7.2.1) en général n'entraîne pas
(5.7.2.2) pour toute générisation $x'$ de $x$, vu qu'on n'a aucune relation
d'inégalité entre $\operatorname{prof}(\mathscr F_x)$ et
$\operatorname{prof}(\mathscr F_{x'})$ (0, 16.4.6).

\item[(iv)] Il résulte aussitôt de la définition que si $\mathscr F$ vérifie
$(S_k)$ en un point $x$, il vérifie aussi $(S_k)$ en tout point $x'$
générisation de $x$.

\item[(v)] La propriété $(S_k)$ est surtout importante pour $k=1$ et
$k=2$; elle a été introduite pour $k=2$ par Serre, pour exprimer son critère
de normalité (cf. (5.8.5)).

\item[(vi)] Soient $X$ un préschéma localement noethérien, $Y$ un
sous-préschéma fermé de $X$, $j:Y\to X$ l'injection canonique, $\mathscr G$
un $\mathcal O_Y$-Module cohérent. Il est clair que pour tout $x\in Y$, on a
$\dim(\mathscr G_x)=\dim((j_*(\mathscr G))_x)$ et
$\operatorname{prof}(\mathscr G_x)=\operatorname{prof}((j_*(\mathscr G))_x)$,
d'où
$\operatorname{coprof}(\mathscr G_x)=\operatorname{coprof}((j_*(\mathscr G))_x)$.
Pour que $\mathscr G$ vérifie $(S_k)$, il faut et il suffit que
$j_*(\mathscr G)$ vérifie $(S_k)$.
\end{enumerate}
\end{remarks}

\begin{proposition}[5.7.4]
\label{IV.2.s5.7.4-fr}
Soient $X$ un préschéma localement noethérien, $\mathscr F$ un
$\mathcal O_X$-Module cohérent. Posons $S=\operatorname{Supp}(\mathscr F)$,
et pour tout $n\geq0$, désignons par $Z_n$ l'ensemble des $z\in X$ tels que
$\operatorname{coprof}(\mathscr F_x)>n$, de sorte que $Z_n\subset S$.
\begin{enumerate}
\item[(i)] Pour que $\mathscr F$ vérifie la propriété $(S_k)$, il faut et il
suffit que, pour tout $n\geq0$, on ait
\begin{equation}
\operatorname{codim}(Z_n,S)>n+k.
\tag{5.7.4.1}\label{IV.2.s5.7.4.1-fr}
\end{equation}
\item[(ii)] Supposons en outre que les $Z_n$ soient fermés dans $X$. Alors,
pour que $\mathscr F$ vérifie la propriété $(S_k)$ en un point $x\in S$, il
faut et il suffit que, pour tout $n\geq0$, on ait
\begin{equation}
\operatorname{codim}_x(Z_n,S)>n+k.
\tag{5.7.4.2}\label{IV.2.s5.7.4.2-fr}
\end{equation}
\end{enumerate}
\end{proposition}

(i) On a en effet, par définition (5.1.3),
$\operatorname{codim}(Z_n,S)=\inf_{z\in Z_n}(\dim(\mathcal O_{S,z}))$ et
l'inégalité (5.7.4.1) signifie donc (5.1.12.2) que, pour tout $z\in X$, et
tout $n\geq0$, la relation
\[
\dim(\mathscr F_z)-\operatorname{prof}(\mathscr F_z)\geq n+1
\]
entraîne la relation
\[
\dim(\mathscr F_z)\geq n+k+1.
\]
Mais si l'on pose $a=\dim(\mathscr F_z)$,
$b=\operatorname{prof}(\mathscr F_z)$, on a $b\leq a$, et dire que pour tout
$n\geq0$, la relation $b\leq a-n-1$ implique $k\leq a-n-1$ équivaut à dire
que $b\geq\inf(k,a)$, d'où la proposition.

